%!TEX program = lualatex
\documentclass[ngerman,headheight=70pt]{scrartcl}
\usepackage[ngerman]{babel}
\usepackage{blindtext}
\usepackage[utf8]{luainputenc}
\usepackage{scrlayer-scrpage}
\usepackage[T1]{fontenc}
\usepackage{fontspec}
\usepackage{multicol}
\usepackage{enumitem}
\usepackage[german=quotes]{csquotes}
\usepackage[bottom=2cm,footskip=8mm]{geometry}
\usepackage{hyperref}
\setmainfont{Source Sans Pro}
\MakeOuterQuote{"}

\setlength{\parindent}{0pt}
\setlength{\parskip}{6pt}

\pagestyle{plain.scrheadings}
\thispagestyle{scrheadings}
\lohead{\vspace*{-1.5cm}\includegraphics[scale=0.4]{../images/stupa_logo_rot.pdf}}
\rohead{\\Präsidium des\\Studierendenparlaments}

\lofoot[Universität Hamburg · Präsidium des Studierendenparlaments\\\hrule Von-Melle-Park 5 · D-20146 Hamburg]{Universität Hamburg · Präsidium des Studierendenparlaments\\\hrule Von-Melle-Park 5 · D-20146 Hamburg}
\cofoot[]{}
\cohead[\thepage]{}

\renewcommand{\thesubsection}{TOP \arabic{subsection}}
\renewcommand{\thesubsubsection}{\arabic{subsubsection}.}
\setcounter{subsection}{-1}

\begin{document}
    UHH · StuPa-Präsidium · Von-Melle-Park 5 · D-20146 Hamburg

    %%% PAGE ONE (EDIT LATER) %%%%
    % Hier den Titel des Protokolls eintragen
    \section*{Protokoll der 6. Sitzung des Studierendenparlaments vom 30. Juni 2016}

    % Protokollanten eintragen
    \textbf{Protokoll: Jim Martens}\\
    \textbf{Ort: S28 in der ehemaligen HWP}\\
    \textbf{Beginn: 18.24 Uhr}\\
    \textbf{Ende: 22.57 Uhr}

    % Wohl die komplizierteste Stelle im Protokoll
    % Als Faustregel sollten nicht mehr als 3 Namen auf einer Zeile stehen
    % Einfach an dieser Vorlage orientieren
    % Im Idealfall ist hier nicht viel zu ändern, da kaum jemand fehlt.
    % Rücktritte müssen natürlich pro Protokoll angepasst werden
    \vspace{0.5cm}
    \begin{tabular}{ll}
        Anwesend: & \\
             RCDS (5 Sitze): & Ramon Weilinger, Jennifer Maack \\
             CampusGrün (14 Sitze): & Laura Franzen, Geoffrey Youett, Elena Rysikova, \\
                                   & Yasemin Günther, Melf Johannsen, Tahnee Herzig, \\
                                   & Mario Moldenhauer, Christian Lagod, Jim Martens,\\
                                   & Svenja Horn, Mirzo Ulugbek Khatamov, Armin Günther \\
             BGZM-Liste (2 Sitze) : & Jakob Pape \\
             WiWi (2 Sitze): &  \\
             Unicorns (5 Sitze): & Johannes Peplow, Annkathrin Löffler \\
             Liste LINKS (3 Sitze): & Gunhild Berdal, Till Petersen, Sinah Mielich \\
             HWP (2 Sitze): &  \\
             MIN (4 Sitze): & Ailina Salten \\
             SDS* (3 Sitze): & Mena Winkler, Jacob Petersein, Artur Brückmann \\
             Bart-LISTE (2 Sitze): & Timo Zeimet \\
             LHG (1 Sitz):  & Tobias Heisig \\
             harte zeiten (1 Sitz): & Tobias Berking \\
             Jura (1 Sitz): & \\
             AL (2 Sitze): & Karima Schulze, Henri Weber \\
            & \\
        Entschuldigt: & Antonia Niecke, Ramin Shakiba (RCDS), Philipp Droll, \\
                      & Martin Sievert (CG), Karen Martirosian (BGZM), Katharina Kucza, \\
                      & Andreas Hartkamp, Marielle Hermstrüwer (Unicorns), Samet Gunay, \\
                      & Ajdina Karahasan (HWP), Nabil Jabareen, Jan Detampel (MIN),\\
                      & Dominic Laumer (Bart), Johann Baumhoefener (Jura) \\
                      &\\
        Unentschuldigt abwesend: & Benjamin Welling (RCDS), Claas-Friso Hente, Kay Zöllmer (WiWi), \\
                                 & Lotte Rullkötter \\
                                 &\\
        Rücktritte: & Thea Wahlers (MIN) \rightarrow Nabil Jabareen\\
    \end{tabular}
    \newpage
    %%% Vorgeschlagene Tagesordnung %%%
    \underline{Vorgeschlagene Tagesordnung}
    \begin{enumerate}[label={\textbf{Top \theenumi}},leftmargin=*]
        % folgende Zeile bei Sitzungsfortsetzungen einkommentieren und x durch
        % die Zahl des zuletzt behandelten TOPs ersetzen
        %\setcounter{enumi}{x}
        \item Diskussion mit Wissenschaftssenatorin Katharina Fegebank (60 Min.)
        \item Haushalt 2016/2017 (V1617-025, V1617-030) (75 Min.)
        \begin{enumerate}
            \item 2. Lesung (70 Min.)
            \item 3. Lesung (5 Min.)
        \end{enumerate}
        \item Nachtragshaushalt 2015/2016 (V1617-026, V1617-030) (25 Min.)
        \begin{enumerate}
            \item 2. Lesung (20 Min.)
            \item 3. Lesung (5 Min.)
        \end{enumerate}
        \item Feministischer Juristinnentag 2017 in Hamburg (V1617-028) (30 Min.)
        \item Weiter initiativ mit der Strategie "Uni für Alle statt Exzellenz" (V1617-029) (30 Min.)
        \item Verlegung der StuPa-Wahl in das Sommersemester (V1617-015) (60 Min.)
        \item Verschiedenes
    \end{enumerate}

    \newpage

    % Hier beginnt Teil A. Bei Fortsetzungen kommt hier nur der Teil über
    % die Feststellung der endgültigen Fassung der Tagesordnung.
    % Für sonstige Sitzungen ist dies bereits das Gerüst für die Formalia.
    % Man muss nur noch die protokollierten Inhalte eintragen.

    {\Large \underline{Teil A}}

    \subsection{Formalia}

    \textit{Gunhild eröffnet die Sitzung um 18.24 Uhr.}

    \subsubsection{Geschäftsbericht Präsidium}

    Das Präsidium hat sich Gedanken über die bevorstehende Diskussion mit
    Frau Fegebank gemacht. Die Exzellenzstrategie und Finanzierung der Hochschulen
    werden in einem Themenblock zusammengefasst. Frau Fegebank wird eingangs etwas
    vortragen, das Ziel ist aber schnell in die Diskussion einzusteigen.

    In den kommenden Wochen hat das Präsidium vor sich Gedanken über die Wahlen
    zu machen. Öffentlichkeitsarbeit soll ebenso Thema sein.

    \subsubsection{Anfragen an das Präsidium}

    Keine Fragen an das Präsidium.

    \subsubsection{Geschäftsbericht AStA}

    Franziska hält den Geschäftsbericht des AStA. Fragen könnten jedoch an einzelne
    Referent*innen gestellt werden.

    Am 23. Juni war die Wissenschaftsausschusssitzung, in der sich der Ausschuss
    mit dem Studierendenwerk beschäftigt hat. Franziska berichtet von der Sitzung,
    in der Moritz und Artur auf der Senatsbank saßen und einige gute Punkte für
    das Studierendenwerk anbringen konnten. Darunter fällt z.B. die Frage
    an Herrn Allemeyer, ob sich das Studierendenwerk bereits im Zusammenhang
    mit Flüchtlingen einbringt. Allemeyer entgegnete, dass dies bislang nicht
    der Fall sei, das Studierendenwerk sich aber gerne einbringen würde. Dafür
    wären jedoch entsprechende Mittel notwendig.

    Folgend berichtet Franziska über das EM-Begleitprogramm, welches gut vorankommt.
    In dem Programm werden die Länder der gegnerischen Mannschaften der
    deutschen Fußballnationalmannschaft in der EM 2016 detailliert dargestellt.
    Es wird ein Fokus auf die soziale und politische Situation in den Ländern
    gelegt. Im Zusammenhang damit gab es auch ein Treffen mit französischen
    Studierenden, die Mitglied einer der großen studentischen Gewerkschaften
    Frankreichs sind.
    Im Rahmen dieses Programms gibt es am 9. Juli um 10 Uhr den "Here to Stay Cup 2016"
    auf den Kunstrasenplätzen hinter dem Millerntorstadium. Dieses Fußballturnier
    richtet sich an Refugee-Mannschaften, studentische Teams und alle die gerne
    Fußball spielen möchten, aber keine Lust auf Rassismus haben.
    Am 14. Juli gibt es um 18 Uhr eine Veranstaltung namens "EM und Nationalismus"
    im Raum 103 im Pferdestall (Allendeplatz 1).

    Der AStA hat die Aktivitäten zur Exzellenzinitiative (neu: Exzellenzstrategie)
    ausgewertet. In dem Zusammenhang gibt es am 7.7. eine Veranstaltung zur
    Exzellenzstrategie namens "Von libertären Mauerkaninchen und totalitären
    Schäferhunden: Wie zur Konformität im Wissenschaftsbetrieb erzogen wird".
    Die Veranstaltung findet um 18 Uhr in Hörsaal G des Philosophenturms statt.

    Franziska fährt den Bericht fort und berichtet über die studentische Arbeit
    gegen Rechts. Im Rahmen dieser Arbeit werde sich an einer Anti-AfD-Veranstaltung
    in Rahlstedt beteiligt und Flugblätter verteilt.
    Am 5. Juli findet der vom Ausschuss gegen Rechts organisierte Abend gegen
    Rechts statt.

    Der gesamte AStA wird sich in den nächsten Wochen zu einer Klausursitzung
    treffen und die Frage "Wie gegen rechts arbeiten" erörtern.

    \textit{Zeitansage: 18.30 Uhr}

    Golnar berichtet aus dem Referat für internationale Studierende. Es hat eine
    Veranstaltung gegeben zur russischen Linken und ihrer Lage im aktuellen Russland
    von Putin. Ein russischer Metallarbeiter war zu Gast. Die Veranstaltung hat
    im Hörsaal der ehemaligen HWP stattgefunden und war gut besucht.
    In der kommenden Zeit werde eine politische Vollversammlung stattfinden.
    Zudem bestehe eine akute Abschiebungsbedrohung für viele ausländische
    Studierende. In dem Zusammenhang gibt es im Juli einen Termin der Harburger
    Ausländerbehörde.

    Geoffrey führt den Rohschnitt eines die Demo gegen die Exzellenzinitiative
    zusammenfassenden Videos vor.

    \subsubsection{Anfragen an den AStA}

    Till fragt, ob sich die Studierenden der französischen Gewerkschaft zum
    Verhältnis von der Gewerkschaft zu der allgemeinen Studierendenbewegung
    geäußert haben und inwiefern solche Äußerungen für das Hamburger Rotstiftbündnis
    angewendet werden können.

    Artur antwortet, dass es in Frankreich nicht an jeder Hochschule politische
    Gruppen gebe, sondern stattdessen nationale Gewerkschaften. Die größte
    dieser Gewerkschaften ist UNEF\footnote{Union nationale des étudiants de France\\
    \url{https://de.wikipedia.org/wiki/Union\_nationale\_des\_\%C3\%A9tudiants\_de\_France}}.
    Studierende werden als intellektuelle Arbeiter begriffen. Die heutige UNEF
    pflegt eine Nähe zu anderen Gewerkschaften. Sie hat sich intensiv mit dem
    neuen Arbeitsgesetz in Frankreich\footnote{u.a. der Spiegel dazu:
    \url{http://spon.de/aeInq}}
    auseinandergesetzt. Mittlerweile hat sich ein breites Bündnis gegen das Gesetz
    gegründet. Einer der französischen Studierenden habe gesagt, dass alle
    Gewerkschaften zwei Demos pro Wochen machen sollten. Dies könne den bereits
    abflauenden Aktivitäten um das Gesetz einen neuen Schub geben.

    Golnar merkt an, dass sich die Aktionsformen der Studierenden in Frankreich
    nicht grundlegend von unseren Aktivitäten unterschieden. Es sei somit eine
    Haltungsfrage. Es würden sich in Frankreich keine Gedanken darüber gemacht
    werden, ob z.B. ein Gesetz die konkrete Klientel etwas angehe. Da es
    Gesellschaftspolitik sei, gehe es somit alle etwas an.
    Golnar äußert auch Kritik an den französischen Studierenden. Sie würden sich
    zu wenig mit Militarismus auseinandersetzen und seien zu sehr aktionistisch
    unterwegs. Außerdem fehle eine Positivorientierung.

    Till fragt, ob es besondere Erkenntnisse bei den Veranstaltungen des
    EM-Begleitprogramms gegeben habe.

    \subsubsection{Dringlichkeitsanträge des AStA}

    Keine.

    \subsubsection{Aktuelle Stunde}

    Es gab keinen Antrag.

    \subsubsection{Feststellung der endgültigen Fassung des Teils B der Tagesordnung}

    Das Präsidium schlägt vor den Top 3 Nachtragsaushalt vorzuziehen auf Top 2.
    Dagegen gibt es kein Widerspruch. Die Kritischen Jurastudierenden haben ihren
    Antrag zurückgezogen. Die so geänderte Tagesordnung wird ohne Widerspruch
    angenommen.

    %%% BESCHLOSSENE TAGESORDNUNG %%%
    \underline{Endgültige Fassung der Tagesordnung}
    \begin{enumerate}[label={\textbf{Top \theenumi}},leftmargin=*]
        % folgende Zeile bei Sitzungsfortsetzungen einkommentieren und x durch
        % die Zahl des zuletzt behandelten TOPs ersetzen
        %\setcounter{enumi}{x}
        \item Diskussion mit Wissenschaftssenatorin Katharina Fegebank (60 Min.)
        \item Nachtragshaushalt 2015/2016 (V1617-026, V1617-030) (25 Min.)
        \begin{enumerate}
            \item 2. Lesung (20 Min.)
            \item 3. Lesung (5 Min.)
        \end{enumerate}
        \item Haushalt 2016/2017 (V1617-025, V1617-030) (75 Min.)
        \begin{enumerate}
            \item 2. Lesung (70 Min.)
            \item 3. Lesung (5 Min.)
        \end{enumerate}
        \item Weiter initiativ mit der Strategie "Uni für Alle statt Exzellenz" (V1617-029) (30 Min.)
        \item Verlegung der StuPa-Wahl in das Sommersemester (V1617-015) (60 Min.)
        \item Verschiedenes
    \end{enumerate}

    \subsubsection{Feststellung der Beschlussfähigkeit}

    \textit{Das StuPa ist mit 28 anwesenden Parlamentarier*innen beschlussfähig.}

    \subsubsection{Genehmigung der Protokolle der vorangegangenen Sitzungen}

    Das Protokoll der vierten Sitzung wird ohne Anmerkungen oder Widerspruch
    genehmigt.

    \textit{Zeitansage: 19.00 Uhr}

    \vspace{0.5cm}
    {\Large \underline{Teil B}}

    % Teil B beginnt
    % Dieser wird sich für jede Sitzung unterscheiden.
    % Jeder TOP ist eine subsection.
    % Vorhandene Beispiele anfangs entfernen und durch tatsächliche TOPs ersetzen

    % Folgende Zeile bei Sitzungsfortsetzungen einkommentieren und x mit
    % der Zahl des zuletzt behandelten TOPs ersetzen.
    %\setcounter{subsection}{x}
    \subsection{Diskussion mit Wissenschaftssenatorin Katharina Fegebank}

    \textit{Frau Fegebank trifft sehr pünktlich drei Minuten vor 19 Uhr ein.}

    Fegebank beginnt mit ein paar einführenden Worten. In dieser Einführung
    stellt sie klar, dass sie niemals gegen die Exzellenzinitiative war, sondern
    nur gegen die konkrete Ausgestaltung. Sie lobt das Verhandlungsergebnis
    und bezeichnet es als Erfolg für die gesamte Hochschullandschaft und den
    Hochschulstandort Hamburg. Die Hamburger Universität könne mit viel
    Selbstbewusstsein in die nächste Runde gehen. Außerdem habe sie bereits
    zwei Exzellenzcluster.

    Desweiteren geht sie auf die Erhöhung der Mittel für die kleineren
    Hamburger Hochschulen ein und spricht von Erhöhungen um 3\% bzw. 4\%.
    Sie geht ebenfalls auf die mögliche destruktive Wirkung der Exzellenzinitiative
    ein und sagt, dass die Forschungsförderung diese Binnenwirkung der
    Exzellenzinitiative verhindern soll.

    Zudem würden Gespräche über einen "Aufwuchs" mit der TUHH und der HAW geführt.
    Schließlich setzt sie sich für die Verstetigung der Hochschulpakt-Mittel ein.

    Gunhild dankt für die einführenden Worte.

    Till führt an, dass es entgegen der Aussage des Hochschulrates durchaus
    Kritik aus der Uni an der Exzellenzinitiative gab. Er wirft die Frage auf,
    warum Exzellenz nötig sei und man nicht für alle bessere Bedingungen schaffe.
    Ferner fragt er, welches Verständnis von wissenschaftlicher Qualität hinter
    der Exzellenzinitiative stecke und woher das Verständnis komme, dass
    Wissenschaft am besten funktioniert, wenn es Gewinner und Verlierer gibt.

    \textit{Zeitansage: 19.15 Uhr}

    Jochen erzählt Einiges zur Berliner TU und ihrer geschichtlichen Vergangenheit
    und den Lehren, die daraus gezogen wurden. Entgegen der Berliner TU sei die
    Hamburger sehr ingenieurwissenschaftlich und bei einem Wachstum der TU sollte
    die gesellschaftliche Bedeutung von Technologie evaluiert werden.

    Franziska erwähnt den sog. "Double Speak": Nach außen wird schöne Antragsprosa
    geschrieben, aber hinter vorgehaltener Hand wird sich stark über die
    Exzellenzinitiative beschwert. Im Imboden-Bericht werde erwähnt, dass
    sich wissenschaftliche Mitarbeiter in ihren produktivsten Jahren mit
    dem Schreiben von Anträgen und dem Jagen von wenigen Festanstellungen
    beschäftigten.

    Laura fragt als letzte Person in der ersten Fragerunde, wie man sich an einer
    Uni mit Exzellenzinitiative und wenigen Festanstellungen eine Familie vorstellen
    solle.

    Fegebank geht in langen Zügen auf die aufgeworfenen Punkte und Fragen
    mehr oder weniger direkt ein. In Bezug zu Till sagt sie, dass sie es nicht
    schlecht finde, dass es Wettbewerb gibt. Im Prinzip resümiert sie diesen
    Austausch mit der Formel "Agree to disagree" und sagt damit effektiv aus,
    dass eine Einigung diesbezüglich zwischen Tills Ansicht und ihrer Ansicht
    ausgeschlossen sei. Insgesamt erwähnt sie häufig den Standort Hamburg.
    Ebenfalls äußert sie sich positiv zur Profilierung von Hochschulen. Als Beispiel
    erwähnt sie, dass eine Hochschule besonders gut in Lehre sein könne.

    Wer sich für eine detailliertere Fassung interessiert, sei an die Aufzeichnung
    der Sitzung\footnote{link} verwiesen. Die Antwort von Frau Fegebank beginnt
    bei Minute xy.

    \textit{Zeitansage: 19.30 Uhr}

    Die zweite Fragerunde beginnt mit Tobias, welcher die Notwendigkeit der
    Exzellenzinitiative in Frage stellt. Insbesondere äußert er, dass das
    Bewerbungsverfahren für die Exzellenzinitiative einem Erziehungsverfahren
    entspreche. Stattdessen müsse Wissenschaft Bestehendes in Frage stellen.
    Provokativ stellt er die Frage, warum es Wettbewerb brauche und nicht
    Kooperation gelebt werde.

    Gunhild führt aus, dass eine höhere Wettbewerbsfähigkeit nicht unbedingt
    auch besser für die Menschheit ist. Ebenfalls kann sie nicht verstehen,
    wie eine Spitzenförderung die Vielfalt fördere. Auch seien Drittmittel ein
    Problem für die Unabhängigkeit von Forschung.

    Artur greift Aspekte von Tobias und Gunhild auf und weist darauf hin, dass
    die Exzellenzinitiative eine "Mainstreamisierung" verursacht. Auch werde
    seitens der Hochschulen hauptsächlich auf die Mittel geachtet, die Inhalte
    seien nachgelagert.

    Golnar erwähnt ein Demokratiedefizit in dem Verfahren, da die Entscheidung
    über die Exzellenz nicht von demokratisch gewählten Vertreter*innen gemacht
    werde. Sie fragt wie Frau Fegebank zu diesem Defizit steht. Golnar greift
    ebenso den Aspekt der Vielfalt auf und fragt explizit wie Vielfalt und
    Exzellenzinitiative zusammenpassten.

    \textit{Zeitansage: 19.46 Uhr}

    Ramon durchbricht die Reihe, indem er sich für die Exzellenzinitiative
    ausspricht. Er fragt, ob der Exzellenzbegriff nicht zu sehr auf
    naturwissenschaftliche Leistungen fixiert sei und ob die Kritik an der
    Exzellenzinitiative nicht zu spät kam.

    Frau Fegebank geht direkt auf diese Frage ein und erklärt, dass von Hamburg
    von Beginn an Kritik geäußert worden sei. Sie erwähnt, dass sie noch keine
    Infos zur Wissenschaftsratsbegehung habe.

    Auf kritische Anmerkungen antwortet sie mit der Frage, was denn passieren würde,
    wenn man der Uni auf einen Schlag 0,5 Milliarden Euro geben würde. Sie stellt
    fest, dass zunächst Party gemacht würde, aber was käme danach? Die Sichtbarkeit
    spiele eine große Rolle bei der Hochschulfinanzierung, da es Leute gebe, die
    hinterfragen, warum so viel Geld in die Hochschulen gegeben wird.

    Ailina fragt Fegebank zu Arbeitssicherheit und den Kosten für Investitionen,
    da es in der Biologie viele kaputte Labore gebe und bislang nichts dagegen
    getan werde.

    Anschließend zieht Till ein Zwischenfazit.

    Fegebank kündigt an, dass sie sich gerne die Biologie anschaue.

    \textit{Zeitansage: 20.08 Uhr}

    \textit{Pause bis 20.20 Uhr, verlängert bis 20.25 Uhr}

    \subsection{Nachtragshaushalt 2015/2016}
    \subsubsection{2. Lesung}

    Ramon erklärt das Prozedere für die Abstimmungen. Melf stellt das erste
    Paket für den Nachtragshaushalt (NHH) vor.

    Keine Gegenrede zum Paket.

    Einstimmig wird der Beschlussempfehlung gefolgt.

    Melf stellt das zweite Paket vor. Keine Gegenrede. Einstimmig wird der
    Beschlussempfehlung gefolgt.

    \subsubsection{3. Lesung}

    M:1:0 wird NHH angenommen.

    \subsection{Haushalt 2016/2017}

    \subsubsection{2. Lesung}

    Melf stellt die Pakete für den Haushalt 16/17 vor. Melf stellt das erste
    Paket für den Haushalt genauer vor. Nico hält die Fürrede für den Antrag
    den Topf für Projektkräfte (PKs) zu erhöhen.

    Tobias hält Fürrede für Antrag den Topf für Referent*innen (Refs) zu erhöhen.
    Er steigt ein in die Debatte.

    \textit{Zeitansage: 20.40 Uhr}

    Karima fragt Till, ob er sich nicht um die Pluralität Sorgen macht, wenn
    7 1/2 PK-Stellen für 2 Ref-Stellen wegfallen.

    Till möchte nicht, dass bei PKs und Refs über Stellen gesprochen wird.
    Er ist der Meinung, dass Geld auf Dauer keine Lösung ist.

    Tobias verweist auf AStA in Lüneburg, der keine PKs hat.

    M:5:4 für Beschlussempfehlung.

    Melf stellt das zweite Paket vor. Er stellt seine Meinung zum Änderungsantrag
    zu Aktionen vor. Er stellt auch den Punkt "Vereinsmitgliedschaften FSRe" vor.

    Tobias hält Fürrede für den Antrag zu Aktionen und Veranstaltungen der VS.

    Melf hält ein Plädoyer für einen realistischen Haushalt, wo nur so viel
    Geld in Posten gesteckt wird, wie benötigt wird.

    Till hält ein Plädoyer für eine Erhöhung des Aktivitätentopfes.

    M:8:2 für Beschlussempfehlung.

    Melf stellt das dritte Paket FZS vor und hält eine Fürrede. Vincent hält
    eine Gegenrede zu den Anträgen und sagt, dass Vollmitglied = Vollmitglied
    bedeuten muss. Erhöhung im Personalbereich aber sinnvoll.

    Geoffrey pflichtet Vincent bei, dass es sinnvoll wäre ein Ref-Stelle zu haben.
    Der Betrag sollte nicht verringert werden. Er geht auch auf das Argument ein,
    dass andere Studierendenschaften weniger zahlen. Diese lasse sich nicht
    eindeutig verifizieren.

    \textit{Zeitansage: 21.00 Uhr}

    Artur geht auf den europäischen Zusammenhang ein.

    Till führt eine Schauspielleistung vor und sagt, dass Geschacher reines
    Marktprinzip und daher abzulehnen ist.

    Nach Protesten gegen Rechte in Sachsen wurden die Studierendenschaften
    in Sachsen mit den Stimmen der NPD verboten und degradiert auf reine
    Vereinsebene.

    6:M:2 gegen Beschlussempfehlung. \\
    7:M:2 gegen Antrag Mitgliedsbeitrag zu senken. \\
    1:14:viele gegen Erhöhung Aufwandsentschädigung.

    Melf stellt das vierte Paket vor. Till bittet in Zukunft Begründungen
    so zu verfassen, dass sie verständlich sind.

    M:0:1 für Beschlussempfehlung.

    Es folgt die finale Abstimmung über den Haushalt, wo die Mehrheit
    bei einer Gegenstimme und vier Enthaltungen für den Haushalt stimmt.

    \textit{Zeitansage: 21.17 Uhr, Pause bis 21.30 Uhr, verlängert bis 21.35 Uhr}

    \subsection{Weiter initiativ mit der Strategie "Uni für Alle statt Exzellenz"}

    \textit{Fraktionspause Liste Links 3 Minuten}

    Till stellt den Antrag vor.

    \textit{Zeitansage: 21.45 Uhr}

    Artur meint, dass es eine weitere Entstellung zur Kenntlichkeit
    gegeben hat.
    Die Differenz zwischen "Trockenschwimmen" und sinnvoller Forschung müsse
    herausgestellt werden. Fegebanks Punkt über die Sichtbarkeit der Uni habe
    einen wahren Kern. Allerdings sollte es in der Sichtbarkeit nicht um
    Sonderforschungsbereiche, sondern um qualitativ sinnvolle Aktionen wie z.B.
    Engagement für Flüchtlinge gehen.
    Man müsse besser klarmachen, dass eine Alternative zur Exzellenzinitiative
    real möglich sei. Man könne von der Bewegung um Bernie Sanders lernen.

    Tobias stellt die historischen Zusammenhänge dar.

    Armin meint, dass man generell das Drittmittelsystem bekämpfen müsse.

    Till befürwortet eine solche allgemeine Kritik. Er geht auf die Diskussion
    mit Fegebank ein. Sie sei deutlich ins Schwimmen geraten. Vor allem
    habe sie mit dem Versprechen von 500 Mio. Euro einen Fehler begangen, den
    man als Senatorin nicht machen sollte. Die Abfolge ihrer Äußerung sei wie folgt:

    An sich ist Mittelerhöhung richtig \(\rightarrow\) Oh das darf ich nicht
    sagen \(\rightarrow\) Habe ich gar nicht erst gemeint.

    Till würde gerne die Diskussion mit Fegebank auswerten.

    Mario unterstützt die Ansicht von Armin und möchte vor der Abstimmung den
    genauen Text wissen.

    Geoffrey findet Fegebanks Beispiel mit den 500 Mio. Euro zynisch. Vor allem
    der Widerspruch, dass die 500 Mio. zu wenig wären, die aktuelle Situation
    im Vergleich dazu aber besser, ärgert ihn. Insbesondere die angestrebte
    Profilbildung macht ihm Sorgen.

    \textit{Zeitansage: 22.05 Uhr}

    Yasemin hat den Eindruck, dass sie sehr unter Druck stand. Sie ist ein
    wenig enttäuscht von dem Besuch. Fegebank habe überhaupt nicht
    gesagt, dass soziale Wissenschaften wichtig seien.

    Jim stellt fest, dass Fegebank mit einem neoliberalem Werbebeitrag
    begonnen habe und langsam in eine echte Diskussion eingestiegen sei.
    Sie schien zudem ein wenig gefangen zu sein zwischen Senatsvertretung und
    persönlicher Ambition.

    Laura äußert sich über die Profiloberstufe und wie desaströs deren Einführung
    gelaufen sei. Ferner sei es interessant, dass Fegebank kein einziges Mal das
    Wort "Effizienz" benutzt habe, obgleich die Aussagen Fegebanks eigentlich
    die ganze Zeit damit zu tun gehabt hätten.

    Armin weist darauf hin, dass Fegebank häufig betont habe, dass in dem Hamburger
    Haushalt eine Erhöhung der Grundfinanzierung um 3\% bis 4\% geplant sei, obwohl
    die tatsächliche Geldsumme lächerlich sei.

    Golnar erwähnt, dass Frau Fegebank eingangs auf Till entgegnete, dass man eben
    zwei unterschiedliche Religionen habe. Sie hätte dann aber herausgefunden,
    dass an der Uni mit echten Argumenten gearbeitet werde. Man dürfe jetzt nicht
    mit dem Kämpfen aufhören, obwohl es kleine Geldgeschenke gibt.

    Till verweist auf die Demo gegen die Exzellenzinitiative und den Zusammenhang
    zwischen Rathaus und Börse/Handelskammer. Fegebank sei eine junge Senatorin
    und könne sich dem Argument nicht komplett verschließen. Außerdem sei sie
    sehr druckanfällig für beide Seiten.

    Till stellt basierend darauf einen Änderungsantrag vor.

    % TODO ÄA hier einfügen

    Artur stellt seinen Änderungsantrag vor.

    % TODO ÄA hier einfügen

    \textit{Zeitansage: 22.30 Uhr}

    Till erklärt den historischen Werdegang von Rap und HipHop.

    Till übernimmt die Änderungen dem Sinne nach und kündigt an den Text entsprechend
    zu überarbeiten. Der Antrag wird anschließend trotz mangelnder finaler Textfassung
    abgestimmt und eine Mehrheit stimmt bei zwei Gegenstimmen und keinen Enthaltungen
    für den Antrag.

    Die finale Textfassung findet sich im Anhang und gilt entsprechend der Abstimmung
    als beschlossen.

    \subsection{Verlegung der StuPa-Wahl in das Sommersemester}

    \textit{Der Antrag wurde zurückgezogen. Damit enfällt der TOP.}

    \subsection{Verschiedenes}

    Ramon fragt, ob es den Wunsch gibt am 14. Juli eine Stupa-Sitzung zu machen.

    Franziska spricht sich für eine Sitzung des Parlaments aus.

    Melf sieht keine Notwendigkeit für eine Sitzung.

    Till spricht sich für eine Sitzung aus. Er sieht keinen Zusammenhang
    zwischen Franziskas Aussage und Melfs Aussage. Man sollte eine Sitzung machen,
    um noch einmal viel vor dem WS zu machen.

    Jim spricht sich grundsätzlich für eine Sitzung aus, sieht aber Probleme
    bei der Beschlussfähigkeit.

    Tobias meint, dass eine solche Sitzung mit Beschlussfähigkeit möglich sei,
    wenn man es darauf anlege.

    Gunhild spricht sich für eine Sitzung aus und findet, dass Personen, die nicht
    kommen können, sich nicht gegen eine Sitzung aussprechen sollten.
    Gunhild appelliert, dass die Ausschüsse wieder mehr arbeiten sollten.

    Melf spricht sich erneut gegen eine Sitzung aus, findet aber einzelne
    Treffen in der VL-freien Zeit sinnvoll.

    Till findet es wichtig, dass man es auf eine Sitzung anlegt.

    Ramon schlägt vor ein Stimmungsbild zu machen. 10:12:1 für eine weitere Sitzung.

    \textit{Ramon schließt um 22:57 die Sitzung.}

\end{document}
